%=============================================================================
% Wave 3, Iteration 4: Materials Science --- MATBG DFD Enhancement
% AGENT: MATERIALS SCIENCE (W3i4 Standalone Deep-Dive)
%
% W3i3 KEY FINDINGS CARRIED FORWARD:
%   - La3Ni2O7 nickelate: mixed d+s_pm pairing, xi_sc = 3.15 nm, FOM_RT = 252 K
%   - Topological SC: MZM zero-energy degeneracy NOT lifted by DFD
%   - MATBG: simultaneous SC+insulator enhancement is DFD-specific prediction
%   - Nb3Sn + YBCO laminate optimal resonator geometry confirmed
%   - Iron pnictide s_pm pairing (W3i2) carried forward
%   - |lambda-1| < 1.56e-9 (authoritative coupling bound)
%   - a_0 = 1.2e-10 m/s^2 (MOND acceleration scale)
%
% W3i4 TASKS:
%   1. MATBG DFD coupling mechanism in flat-band systems
%   2. Parameter space map (theta, D, n) with DFD phase boundary shifts
%   3. Quantitative T_c enhancement: delta(T_c)/T_c at |lambda-1| bound
%   4. Insulator gap modification: Mott gap change in meV
%   5. Smoking-gun experimental test design
%   6. Flat-band system comparison table (MATBG vs twisted WSe2 vs pentalayer)
%   7. Connection to P2 resonators and P3 quantum computing
%   8. Optimal device geometry for DFD detection
%
% Date: March 2026
%=============================================================================

\documentclass[11pt]{article}
\usepackage{amsmath,amssymb,amsthm}
\usepackage{geometry}
\geometry{margin=1in}
\usepackage{booktabs}
\usepackage{longtable}
\usepackage{array}
\usepackage{enumitem}
\usepackage{hyperref}
\usepackage{xcolor}
\usepackage{bm}
\usepackage{multirow}

\newtheorem{theorem}{Theorem}
\newtheorem{proposition}{Proposition}
\newtheorem{lemma}{Lemma}
\newtheorem{finding}{Finding}
\newtheorem{prediction}{DFD Prediction}
\newtheorem{correction}{Correction}

\newcommand{\ee}[1]{\times 10^{#1}}
\newcommand{\psifield}{\psi}
\newcommand{\avec}{\mathbf{a}}
\newcommand{\kB}{k_{\rm B}}
\newcommand{\Tc}{T_{\rm c}}
\newcommand{\order}[1]{\mathcal{O}(#1)}
\newcommand{\abs}[1]{\left|#1\right|}
\newcommand{\lambdaL}{\lambda_{\rm L}}
\newcommand{\tmagic}{\theta_{\rm magic}}
\newcommand{\nuFill}{\nu}
\newcommand{\DOS}{D(\varepsilon_F)}
\newcommand{\Vmott}{\Delta_{\rm Mott}}
\newcommand{\VDFD}{\delta V_{\rm DFD}}
\newcommand{\lambdabound}{1.56\ee{-9}}
\newcommand{\psizero}{\psi_0}

\title{Wave 3 Iteration 4: Materials Science --- MATBG DFD Enhancement\\
\large Magic-Angle Twisted Bilayer Graphene in the Psi-Field:\\
Parameter Space Mapping, Quantitative Predictions,\\
Smoking-Gun Experimental Tests, and Cross-Patent Integration}
\author{DFD Research Programme --- Agent 17, Materials Science (W3i4)}
\date{March 2026}

\begin{document}
\maketitle
\tableofcontents
\newpage

%=============================================================================
\section{W3i4 Overview and Motivation}
\label{sec:overview}
%=============================================================================

This document is the Wave~3, Iteration~4 dedicated materials science
analysis. The central focus is magic-angle twisted bilayer graphene (MATBG)
as a DFD-sensitive platform: its ultra-flat bands produce an amplified
density of states (DoS) that couples directly to the psi-field, making it
among the highest-sensitivity non-resonator probes of DFD physics. W3i3
identified the simultaneous superconductor (SC) + correlated-insulator
enhancement as a DFD-specific prediction; this document derives that
prediction quantitatively, maps the full parameter space, and designs the
smoking-gun experimental test.

\subsection{Key W3i3 Results Carried Forward}

\begin{itemize}
  \item \textbf{MATBG DoS amplification}: $\sim 2700\times$ above normal
    graphene at $\tmagic \approx 1.1^\circ$, arising from flat-band
    Van~Hove singularity.
  \item \textbf{DFD-specific prediction}: simultaneous enhancement of SC
    and correlated-insulator gaps, which is impossible in standard BCS or
    Hubbard models.
  \item \textbf{Authoritative coupling bound}: $|\lambda-1| < \lambdabound$
    (from Casimir, $g-2$, and atomic spectroscopy).
  \item \textbf{P2 resonator optimal material}: Nb$_3$Sn + YBCO laminate;
    MATBG operates as a \emph{psi-sensor element}, not structural resonator.
  \item \textbf{Nickelate FOM}: La$_3$Ni$_2$O$_7$ at FOM$_{\rm RT} = 252$~K
    remains the highest-$T_c$ DFD candidate for structural SRF.
\end{itemize}

%=============================================================================
\section{DFD Coupling Mechanism in Flat-Band Systems}
\label{sec:mechanism}
%=============================================================================

\subsection{Why Flat Bands Amplify DFD Coupling}

In DFD, the psi-field couples to the local electromagnetic energy density
through the permittivity/permeability modification:
\begin{equation}
  \varepsilon(\mathbf{r}) = \varepsilon_0\bigl[1 + (\lambda-1)\psifield(\mathbf{r})\bigr],
  \quad
  \mu(\mathbf{r}) = \mu_0\bigl[1 + (\lambda-1)\psifield(\mathbf{r})\bigr].
  \label{eq:psi_eps_mu}
\end{equation}
The resulting effective interaction potential between electrons (mediated by
virtual photon exchange through the modified vacuum) gains a psi-field
correction:
\begin{equation}
  V_{\rm eff}(q) = V_{\rm Coulomb}(q)\Bigl[1 + 2(\lambda-1)\psifield\Bigr]
                    + \VDFD(q),
  \label{eq:Veff}
\end{equation}
where the DFD correction to the electron--electron interaction is
\begin{equation}
  \VDFD = (\lambda-1)\,\psifield \cdot \frac{e^2}{\varepsilon_0 q^2}
         \cdot n_{\rm local}.
  \label{eq:VDFD}
\end{equation}
Here $n_{\rm local}$ is the local carrier density. The key amplification
arises because in a flat-band system, the \emph{kinetic energy is
quenched}: $E_{\rm kinetic} \ll U_{\rm Coulomb}$, so interaction effects
dominate and the DFD correction $\VDFD$ is not suppressed by the large
kinetic-energy denominator that operates in normal metals.

\subsection{Flat-Band DoS Enhancement Factor}

In conventional graphene, the DoS at the Fermi level is
$D_0(\varepsilon_F) \approx 4.5\ee{12}$~eV$^{-1}$cm$^{-2}$ at typical
carrier density $n \sim 10^{12}$~cm$^{-2}$. At the magic angle, the flat
band compresses the bandwidth $W_{\rm flat} \approx 5$--$10$~meV (compared
to $W_{\rm Dirac} \approx 1$~eV), giving a DoS enhancement:
\begin{equation}
  \frac{D_{\rm flat}}{D_0} \approx \frac{W_{\rm Dirac}}{W_{\rm flat}}
    \approx \frac{1\,\text{eV}}{7\,\text{meV}} \approx 143.
  \label{eq:DOS_enhancement}
\end{equation}
Including the 8-fold degeneracy (spin $\times$ valley $\times$ sublattice)
and the Van~Hove singularity peak structure, the effective enhancement
relevant for pairing reaches $\sim 2700\times$ (W3i3 result).

\subsection{Process A Sourcing in MATBG}

The psi-field background $\psi_0$ is sourced by Process~A (EM$\to\psi$
conversion) from the ambient EM environment. In MATBG at $T \sim 1$~K, the
dominant EM source is thermal blackbody radiation at the device temperature.
The local psi-field amplitude in a dilution refrigerator environment
(essentially zero thermal photons at GHz frequencies) is governed by the
external laboratory fields and the device's own zero-point fluctuations.
For the purposes of the DFD enhancement calculation, we treat $\psi_0$ as
a dimensionless parameter of order unity normalised to the local metric
perturbation, with the coupling strength set by $|\lambda-1|$.

%=============================================================================
\section{DFD Coupling in Superconducting and Insulating Phases}
\label{sec:dual_phase}
%=============================================================================

\subsection{The Apparent Paradox: Simultaneous SC + Insulator Enhancement}

Standard BCS theory and Hubbard-model physics predict that superconductivity
and Mott insulation are \emph{competing} phases: increasing the attractive
interaction $|V|$ drives Cooper pairing while suppressing Mott physics, and
vice versa. In MATBG, the two phases are observed to coexist in adjacent
regions of the phase diagram, separated by sharp boundaries at integer
filling factors $\nuFill = \pm 2$.

\begin{finding}
  \label{find:simultaneous}
  DFD predicts \emph{simultaneous} enhancement of both the SC gap
  $\Delta_{\rm SC}$ and the correlated-insulator gap $\Vmott$. This is
  not paradoxical because the DFD psi-field acts as a \emph{local
  symmetry-breaking field}: it does not uniformly increase or decrease
  the pairing interaction, but instead modifies the \emph{spatial pattern}
  of the electron density, sharpening the distinction between paired and
  localised regions.
\end{finding}

\subsection{DFD Mechanism for Dual Enhancement}

The psi-field modifies the local electron density through the electrochemical
potential:
\begin{equation}
  \mu_{\rm eff}(\mathbf{r}) = \mu_0 + (\lambda-1)\psifield(\mathbf{r})
    \cdot \frac{\partial E}{\partial n},
  \label{eq:mu_eff}
\end{equation}
where $E$ is the total energy per electron. This position-dependent
chemical potential acts as a spatial modulation of the filling factor.
Near $\nuFill = 2$ (half-filling of the flat bands):
\begin{itemize}
  \item In regions where $\mu_{\rm eff}$ fluctuates \emph{below} the
    Mott gap, electrons become more strongly localised $\Rightarrow$
    insulator gap \emph{enhanced}.
  \item In regions where $\mu_{\rm eff}$ fluctuates \emph{above} the
    Mott gap, electrons are delocalized into SC puddles with enhanced
    pairing $\Rightarrow$ SC gap \emph{enhanced} at phase boundary.
\end{itemize}
This is a spatial phase-separation enhancement: DFD sharpens the
first-order-like transition between SC and insulator phases, increasing
the amplitude of both order parameters at their respective locations.

\begin{prediction}
  \label{pred:dual}
  In dual-gated MATBG, DFD predicts that the SC gap $\Delta_{\rm SC}$
  and the correlated-insulator gap $\Vmott$ both increase by the same
  fractional amount $\delta\Delta/\Delta \approx |\lambda-1|\psi_0 \cdot
  (D_{\rm flat}/D_0)$, where $D_{\rm flat}/D_0 \approx 143$ is the
  flat-band DoS enhancement.
\end{prediction}

%=============================================================================
\section{Quantitative DFD Predictions}
\label{sec:quantitative}
%=============================================================================

\subsection{DFD Correction to the Effective Pairing Potential}

The attractive phonon-mediated pairing interaction in MATBG has been
measured to be $V_{\rm phonon} \approx 3$--$5$~meV (inferred from
$T_c \approx 1$--$3$~K via weak-coupling BCS: $\kB T_c \approx 1.13
\hbar\omega_D e^{-1/\lambda_{\rm BCS}}$, with $\lambda_{\rm BCS} \sim 0.2$).
We adopt $V_{\rm phonon} = 4$~meV as the reference value.

The DFD correction (eq.~\ref{eq:VDFD}) at the flat-band carrier density
$n_{\rm local} \approx 3\ee{12}$~cm$^{-2}$ (half-filling, $\nuFill = 2$)
and using the authoritative bound $|\lambda-1| = \lambdabound$:
\begin{align}
  \VDFD &= |\lambda-1| \cdot \psi_0 \cdot V_{\rm Coulomb}(q \sim k_F)
            \cdot n_{\rm local} \nonumber\\
        &= \lambdabound \times 1 \times
           \frac{e^2}{\varepsilon_0 k_F} \times n_{\rm local}.
  \label{eq:VDFD_estimate}
\end{align}
For MATBG on hBN, the dielectric screening gives an effective
$\varepsilon_r \approx 4$. The Fermi wavevector at half-filling:
$k_F \approx \sqrt{4\pi n / g_s} \approx 1.4\ee{8}$~m$^{-1}$ (where
$g_s = 4$ is the spin-valley degeneracy). Thus:
\begin{align}
  V_{\rm Coulomb}(k_F) &= \frac{e^2}{4\pi\varepsilon_0\varepsilon_r k_F}
    \approx \frac{1.44\,\text{eV\AA}}{4 \times 1.4\ee{-2}\,\text{\AA}}
    \approx 25.7\,\text{meV}.
  \label{eq:VCoulomb}
\end{align}
Therefore:
\begin{align}
  \VDFD &\approx 1.56\ee{-9} \times 25.7\,\text{meV}
    \approx 4.01\ee{-8}\,\text{meV}.
  \label{eq:VDFD_value}
\end{align}
Without flat-band amplification (standard graphene), this is undetectable.
With flat-band DoS enhancement $\mathcal{A} \approx 143$:
\begin{equation}
  \VDFD^{\rm (flat)} = \mathcal{A} \times \VDFD
    \approx 143 \times 4.01\ee{-8}\,\text{meV}
    \approx 5.7\ee{-6}\,\text{meV}.
  \label{eq:VDFD_flat}
\end{equation}

\subsection{T$_c$ Enhancement}

Using weak-coupling BCS, the fractional $T_c$ enhancement is:
\begin{equation}
  \frac{\delta T_c}{T_c} = \frac{\delta V_{\rm DFD}^{\rm (flat)}}
    {V_{\rm phonon} \cdot \lambda_{\rm BCS}^2}
    \approx \frac{5.7\ee{-6}\,\text{meV}}{4\,\text{meV} \times 0.04}
    \approx 3.6\ee{-5}.
  \label{eq:dTc}
\end{equation}
For $T_c = 1.5$~K (typical MATBG), this gives:
\begin{equation}
  \delta T_c \approx 3.6\ee{-5} \times 1.5\,\text{K}
    \approx 5.4\ee{-5}\,\text{K} = 54\,\mu\text{K}.
  \label{eq:dTc_value}
\end{equation}

\begin{finding}
  \label{find:Tc_shift}
  The DFD-induced $T_c$ shift in MATBG at the current bound
  $|\lambda-1| < \lambdabound$ is $\delta T_c \approx 54\,\mu$K.
  State-of-the-art dilution refrigerator thermometry achieves
  $\sim 1$~$\mu$K stability; however, $T_c$ determination in mesoscopic
  MATBG devices is limited by disorder-induced inhomogeneity at the
  $\sim 50$--$100$~mK level. The DFD signal is therefore \emph{below
  current experimental resolution} by a factor of $\sim 1000$.
\end{finding}

\subsection{Mott Gap Enhancement}

The correlated insulator gap at $\nuFill = 2$ is $\Vmott \approx 0.3$--$1$~meV
(measured by Cao et al.\ 2018, Lu et al.\ 2019). The DFD modification of the
Mott gap arises through the effective on-site Hubbard-$U$ modification:
\begin{equation}
  \delta U_{\rm DFD} = |\lambda-1| \cdot \psi_0 \cdot
    \frac{\partial^2 E}{\partial n^2}\Bigg|_{\nuFill=2} \cdot \mathcal{A}.
  \label{eq:deltaU}
\end{equation}
The second derivative $\partial^2 E/\partial n^2 \approx U \sim 10$~meV
(the bare Hubbard $U$). Therefore:
\begin{align}
  \delta U_{\rm DFD} &\approx \lambdabound \times 10\,\text{meV} \times 143
    \approx 2.2\ee{-6}\,\text{meV}.
  \label{eq:deltaU_value}
\end{align}
The fractional Mott gap change:
\begin{equation}
  \frac{\delta\Vmott}{\Vmott} = \frac{\delta U_{\rm DFD}}{U}
    \approx \frac{2.2\ee{-6}\,\text{meV}}{10\,\text{meV}}
    \approx 2.2\ee{-7}.
  \label{eq:dMott}
\end{equation}
For $\Vmott = 0.5$~meV, this gives $\delta\Vmott \approx 0.1$~neV,
far below any accessible measurement.

\begin{finding}
  \label{find:mott_gap}
  The DFD-induced Mott gap change in MATBG is
  $\delta\Vmott \approx 0.1$~neV at $|\lambda-1| = \lambdabound$.
  Tunnelling spectroscopy achieves $\sim 10\,\mu$eV energy resolution;
  the DFD signal is $\sim 10^8$ times smaller. Transport gap measurement
  achieves $\sim 0.1$~meV resolution, still $\sim 10^6$ times larger than
  the signal. \emph{The Mott gap is not a viable DFD detection channel
  in MATBG at the current coupling bound.}
\end{finding}

%=============================================================================
\section{Phase Diagram Parameter Space: $(\theta, D, n)$ Mapping}
\label{sec:phase_diagram}
%=============================================================================

\subsection{Standard MATBG Phase Diagram}

The MATBG phase diagram in the $(n, \theta)$ plane (at fixed $D = 0$) shows:
\begin{itemize}
  \item Correlated insulator at $\nuFill = \pm 2, \pm 3$ (integer flat-band
    filling): Mott-like gaps of $0.3$--$1$~meV.
  \item SC domes centred at $\nuFill \approx \pm 2.2$--$2.7$ with
    $T_c^{\rm max} \approx 1$--$3$~K; dome half-width $\Delta n \approx
    1\ee{12}$~cm$^{-2}$.
  \item SC dome sharpest (and $T_c$ highest) at $\theta = \tmagic
    \approx 1.05^\circ$--$1.15^\circ$.
  \item Displacement field $D \ne 0$ breaks $C_2T$ symmetry and can
    drive the system toward spin-polarised or valley-Hall states.
\end{itemize}

\subsection{DFD Phase Boundary Shifts}

The DFD psi-field shifts the effective chemical potential by
$\delta\mu = (\lambda-1)\psi_0 \cdot \partial E/\partial n$, moving
the SC dome boundary in carrier density by:
\begin{equation}
  \delta n_{\rm SC}^{\rm boundary} = \frac{\delta\mu}{\partial\mu/\partial n}
    = |\lambda-1| \cdot \psi_0 \cdot \frac{\partial E/\partial n}
      {\partial^2 E/\partial n^2} \cdot \mathcal{A}.
  \label{eq:dn_boundary}
\end{equation}
With $\partial E/\partial n \approx \varepsilon_F \approx 5$~meV and
$\partial^2 E/\partial n^2 \approx U \approx 10$~meV:
\begin{align}
  \delta n_{\rm SC}^{\rm boundary}
    &\approx \lambdabound \times 1 \times \frac{5\,\text{meV}}
      {10\,\text{meV}} \times 143 \nonumber\\
    &\approx 1.12\ee{-7}\,\text{(dimensionless fraction)}.
  \label{eq:dn_value}
\end{align}
In absolute carrier density units
(full filling $n_s = 4/A_{\rm moir\acute{e}} \approx 4/(2.5\ee{-10})^2
\approx 6.4\ee{12}$~cm$^{-2}$):
\begin{equation}
  \delta n_{\rm SC}^{\rm boundary}
    \approx 1.12\ee{-7} \times 6.4\ee{12}\,\text{cm}^{-2}
    \approx 7.2\ee{5}\,\text{cm}^{-2}.
  \label{eq:dn_abs}
\end{equation}

\begin{finding}
  \label{find:phase_boundary}
  The DFD shift of the SC dome boundary in MATBG is
  $\delta n \approx 7.2\ee{5}$~cm$^{-2}$. State-of-the-art dual-gated
  devices achieve density control at the $\sim 10^9$~cm$^{-2}$ level
  (gate voltage resolution $\sim 1$~mV). The DFD phase boundary shift
  is $\sim 1000$ times smaller than the measurement resolution.
\end{finding}

\subsubsection{Twist-Angle Shift}

The DFD modification of $T_c$ as a function of twist angle enters through
the bandwidth $W(\theta)$. Near the magic angle:
$W(\theta) \approx W_0 |\theta - \tmagic|/\tmagic$, so
$D_{\rm flat}(\theta) \propto 1/|\theta - \tmagic|$. The DFD-enhanced
$T_c(\theta)$ peak is shifted by:
\begin{equation}
  \delta\theta_{\rm DFD} \approx \frac{\delta T_c}{T_c}
    \cdot \frac{\tmagic}{\partial\ln T_c/\partial\ln W}
    \approx 3.6\ee{-5} \times 0.018^\circ \times 1
    \approx 6.5\ee{-7}\,\text{degrees}.
  \label{eq:dtheta}
\end{equation}
State-of-the-art twist-angle control achieves $\pm 0.01^\circ$; the DFD
angular shift is $\sim 15\,000$ times smaller.

\subsubsection{Displacement Field Shift}

The displacement field $D$ tunes the inter-layer asymmetry. DFD modifies
the effective interlayer hopping $t_\perp \to t_\perp(1 + |\lambda-1|\psi_0)$
by a fractional amount $\lambdabound$. This produces a shift:
\begin{equation}
  \delta D_{\rm DFD} = |\lambda-1| \cdot D_{\rm critical}
    \approx \lambdabound \times 0.5\,\text{V/nm}
    \approx 7.8\ee{-10}\,\text{V/nm}.
  \label{eq:dD}
\end{equation}
Gate voltage control achieves $\sim 1$~mV resolution,
corresponding to $\sim 10^{-4}$~V/nm displacement field resolution.
The DFD shift $\delta D \approx 8\ee{-10}$~V/nm is $\sim 10^5$ times
below experimental reach.

\subsection{Summary Phase Diagram Table}

\begin{table}[h!]
\centering
\caption{DFD phase boundary shifts in MATBG parameter space vs.\ current
experimental resolution. All shifts computed at $|\lambda-1| = \lambdabound$.}
\label{tab:phase_shifts}
\begin{tabular}{lrrrl}
\toprule
Parameter & DFD Shift & Expt.\ Resolution & Ratio & Verdict \\
\midrule
$T_c$ & $54\,\mu$K & $\sim 50$~mK & $10^{-3}$ & Not detectable \\
$\delta n_{\rm SC}$ & $7.2\ee{5}$~cm$^{-2}$ & $10^9$~cm$^{-2}$ & $10^{-3}$ & Not detectable \\
$\delta\Vmott$ & $0.1$~neV & $10\,\mu$eV & $10^{-8}$ & Not detectable \\
$\delta\theta$ & $6.5\ee{-7}$~deg & $0.01^\circ$ & $6\ee{-5}$ & Not detectable \\
$\delta D$ & $8\ee{-10}$~V/nm & $10^{-4}$~V/nm & $10^{-5}$ & Not detectable \\
\bottomrule
\end{tabular}
\end{table}

%=============================================================================
\section{Smoking-Gun Experimental Test Design}
\label{sec:smoking_gun}
%=============================================================================

\subsection{Optimal Strategy: DFD as a Transducer}

Given that direct observation of phase boundary shifts is infeasible at
the current $|\lambda-1|$ bound, the preferred experimental strategy
repurposes MATBG as a \emph{transducer}: a device whose resistance
$R(\nuFill, T)$ near the SC/insulator phase boundary is maximally
sensitive to small changes in the effective pairing potential.

\subsubsection{Near-Critical Resistance Amplification}

At the SC/insulator quantum critical point (QCP) at $\nuFill = 2$,
the resistance exhibits power-law scaling:
\begin{equation}
  R \propto |n - n_c|^{-\nu z},
  \label{eq:R_scaling}
\end{equation}
where $\nu z \approx 1.3$ for the 2D XY universality class. Near $n_c$,
a tiny density shift $\delta n$ produces a resistance change:
\begin{equation}
  \frac{\delta R}{R} = \nu z \cdot \frac{\delta n}{|n - n_c|}.
  \label{eq:dR}
\end{equation}
Biasing at $|n - n_c| \sim \delta n_{\rm DFD} \times \kappa$ with
$\kappa \ll 1$ amplifies the signal, but requires knowing $n_c$ to
sub-$\delta n_{\rm DFD}$ precision.

\subsubsection{Modulated-Psi Protocol}

\begin{prediction}
  \label{pred:mod_psi}
  If an external EM field (Process~A source) modulates the psi-field at
  frequency $f_{\rm mod}$, the MATBG resistance near the QCP should
  exhibit a lock-in response at $f_{\rm mod}$ with amplitude:
  \begin{equation}
    \delta R(f_{\rm mod}) \approx R_N \cdot \nu z \cdot
      \frac{\delta n_{\rm DFD}}{|n - n_c|}
      \approx R_N \times 1.3 \times
      \frac{7.2\ee{5}\,\text{cm}^{-2}}{|n - n_c|},
    \label{eq:dR_lockin}
  \end{equation}
  where $R_N \sim h/e^2 \approx 26\,\text{k}\Omega$ is the resistance
  quantum. Biasing at $|n-n_c| = 10^8$~cm$^{-2}$ (achievable)
  gives $\delta R \approx 26\,\text{k}\Omega \times 9.4\ee{-3}
  \approx 240\,\Omega$. This is measurable in principle but requires
  $\delta n_{\rm DFD}$ to be free of competing noise sources (charge
  noise, gate voltage noise).
\end{prediction}

\subsection{Smoking-Gun: Differential Phase Diagram Comparison}

The true smoking-gun experiment is a \emph{differential comparison}:

\begin{enumerate}
  \item Fabricate two nominally identical MATBG devices at the same
    twist angle $\theta = \tmagic \pm 0.01^\circ$ on hBN.
  \item Place one device inside a mu-metal-shielded enclosure (minimising
    ambient EM fields $\Rightarrow$ reduced Process~A sourcing
    $\Rightarrow$ reduced $\psi_0$).
  \item Place the second device in the unshielded configuration.
  \item Map the full phase diagram $(n, D)$ at $T = 20$~mK for both.
  \item \textbf{DFD prediction}: SC dome boundary differs between shielded
    and unshielded devices by $\delta n_{\rm SC}$ as calculated above.
  \item \textbf{Standard prediction}: no difference (disorder fluctuations
    between devices are indistinguishable from the DFD signal at current
    sensitivity).
\end{enumerate}

\begin{finding}
  \label{find:smoking_gun}
  The differential phase diagram measurement is the most stringent possible
  DFD test in MATBG, but the predicted DFD signal ($\delta n \approx
  7\ee{5}$~cm$^{-2}$) is 3--4 orders of magnitude below the limiting
  noise floor from intrinsic disorder ($\delta n_{\rm disorder} \sim
  10^9$--$10^{10}$~cm$^{-2}$). The experiment can set an upper limit on
  $|\lambda-1|$ only if disorder is reduced by $\sim 10^3$--$10^4\times$
  beyond current state of the art.
\end{finding}

\subsection{Required Material Quality for DFD Detection in MATBG}

To achieve sensitivity to the DFD signal, disorder-limited density
resolution must be reduced to $\delta n_{\rm disorder} < \delta n_{\rm DFD}
\approx 10^6$~cm$^{-2}$. This requires:
\begin{itemize}
  \item hBN encapsulation with atomically flat surfaces (already achieved
    in best devices).
  \item Residual charge inhomogeneity $\delta n < 10^6$~cm$^{-2}$, which
    is $\sim 1000\times$ better than the best current MATBG devices
    ($\delta n \sim 10^9$~cm$^{-2}$).
  \item Gate voltage noise $< 1\,\mu$V (currently $\sim 10$--$100\,\mu$V
    from Johnson noise and 1/f noise at mK temperatures).
\end{itemize}
These requirements lie beyond near-term experimental capability.

%=============================================================================
\section{Optimal Device Design for DFD Detection}
\label{sec:device_design}
%=============================================================================

\subsection{Twist Angle Precision Requirements}

The magic-angle window is $\theta = 1.05^\circ$--$1.15^\circ$ (FWHM
$\approx 0.1^\circ$). To maximise flat-band DoS enhancement, the device
must be within $0.02^\circ$ of the optimal angle. The DFD angular shift
$\delta\theta_{\rm DFD} \approx 6.5\ee{-7}$ degrees does not impose
additional precision requirements (it is far below the fabrication limit).
Required precision: $\pm 0.02^\circ$ (current state of the art with the
``tear-and-stack'' technique achieves $\pm 0.1^\circ$; with recently
developed mechanical rotation stages: $\pm 0.01^\circ$).

\subsection{Substrate and Encapsulation}

\begin{itemize}
  \item \textbf{hBN encapsulation} (top and bottom): reduces charge
    disorder from $\sim 10^{11}$~cm$^{-2}$ (SiO$_2$) to
    $\sim 10^9$~cm$^{-2}$. Required for any DFD-relevant experiment.
  \item \textbf{One-dimensional edge contacts}: avoids disorder at
    graphene/contact interface.
  \item \textbf{Ultra-thin hBN} (2--3 layers, $\sim 0.7$~nm) as dielectric
    spacer between graphite gates: maximises gate capacitance and reduces
    gate voltage noise contribution to density noise.
  \item \textbf{Dual graphite gate geometry}: independent control of
    carrier density $n$ and displacement field $D$ at sub-1\% cross-talk.
\end{itemize}

\subsection{Gate Geometry for Independent $n$ and $D$ Control}

The top gate (TG) and bottom gate (BG) voltages control:
\begin{align}
  n &= \frac{C_{\rm TG} V_{\rm TG} + C_{\rm BG} V_{\rm BG}}{e},
  \label{eq:n_control}\\
  D &= \frac{C_{\rm TG} V_{\rm TG} - C_{\rm BG} V_{\rm BG}}{2\varepsilon_0}.
  \label{eq:D_control}
\end{align}
With matched capacitances $C_{\rm TG} = C_{\rm BG}$, any $(n, D)$ point
in the phase diagram is independently accessible. Optimal design:
$C \approx 30$~nF/cm$^2$ (3-layer hBN dielectric), giving
$\partial n/\partial V \approx 1.9\ee{12}$~cm$^{-2}$V$^{-1}$.

\subsection{Measurement Protocol}

\begin{enumerate}
  \item Cool to $T < 50$~mK in dilution refrigerator.
  \item Apply perpendicular magnetic field $B < 0.1$~T to suppress
    superconducting fluctuations during density mapping.
  \item Map resistance $R(V_{\rm TG}, V_{\rm BG})$ at fixed $T$ to obtain
    full $(n, D)$ phase diagram.
  \item Identify SC dome boundaries at $R < 0.01 R_N$ and insulator
    peaks at $R > 100 R_N$.
  \item Apply external RF excitation at $f_{\rm mod} = 1$--$10$~GHz
    (below the SC gap frequency $2\Delta/h \approx 10$~GHz) to modulate
    ambient EM environment; measure lock-in response of resistance.
  \item Repeat with and without mu-metal shielding.
\end{enumerate}

\subsection{Measurement Discriminator: DFD vs.\ Standard MATBG Physics}

The key measurement that would discriminate DFD from standard MATBG physics
is the \textbf{EM-environment dependence of the phase diagram}:
\begin{itemize}
  \item Standard physics predicts zero dependence on external EM environment
    (at fixed $T$, $n$, $D$, $B$, shielded from direct microwave heating).
  \item DFD predicts a shift in SC dome boundary by $\delta n \approx
    7\ee{5}$~cm$^{-2}$ correlated with ambient EM power via Process~A.
\end{itemize}
This is in principle a \emph{qualitative} smoking gun, because no standard
mechanism produces an EM-dependent SC boundary shift independent of sample
heating. The challenge is eliminating RF heating as a confounder (standard
physics also shifts the boundary if $T$ rises due to absorbed RF power).

\textbf{Control requirement}: demonstrate that the electron temperature
$T_e$ (measured via Coulomb blockade thermometry or Johnson noise
thermometry) does not change when the EM environment is modulated.
Only if $T_e$ is stable and the SC boundary shifts is the DFD explanation
viable.

%=============================================================================
\section{Flat-Band System Comparison Table}
\label{sec:comparison}
%=============================================================================

\begin{table}[h!]
\centering
\caption{Comparison of flat-band moiré systems for DFD signal detection.
FOM$_{\rm DFD}$ is proportional to $\mathcal{A} \times T_c / T_{\rm min}$
where $T_{\rm min}$ is the operating temperature.
All values at current bound $|\lambda-1| < \lambdabound$.}
\label{tab:flatband_comparison}
\renewcommand{\arraystretch}{1.3}
\begin{tabular}{p{3.8cm}rrrrl}
\toprule
System & $T_c$ & $\mathcal{A}$ & $\delta T_c$ & Disorder & DFD \\
       & (K) & & ($\mu$K) & floor & FOM \\
\midrule
MATBG ($\theta=1.1^\circ$) & 1.5 & 143 & 54 & High & Moderate \\
Twisted WSe$_2$ & 0.10 & $\sim 200$ & 4 & Low & Low \\
Pentalayer graphene/hBN & 0.30 & $\sim 80$ & 9 & Low & Low \\
TBG/MoS$_2$ heterostructure & 0.50 & $\sim 150$ & 19 & Medium & Moderate \\
Twisted MoTe$_2$ & 0.07 & $\sim 400$ & 5 & Very low & Low \\
Rhombohedral trilayer/hBN & 0.20 & $\sim 120$ & 10 & Low & Low \\
La$_3$Ni$_2$O$_7$ (nickelate) & 80 & $\sim 3$ & 9000 & Medium & \textbf{High} \\
\bottomrule
\end{tabular}
\end{table}

\subsection{Analysis: MATBG vs.\ Other Systems}

\subsubsection{Twisted WSe$_2$}

Twisted WSe$_2$ at $\theta \approx 5^\circ$ exhibits superconductivity at
$T_c \approx 100$~mK. The flat-band DoS enhancement $\mathcal{A} \approx 200$
exceeds MATBG due to stronger spin-orbit coupling and heavier effective mass.
However, the very low $T_c$ means the system must be operated at
$T < 50$~mK, and the SC gap $\Delta \approx 10\,\mu$eV is comparable to the
DFD-induced gap modification $\delta\Delta \approx 0.01$~neV, giving a
ratio $\delta\Delta/\Delta \approx 10^{-9}$.

\textbf{Verdict}: Twisted WSe$_2$ is a cleaner platform (lower disorder)
but offers no improvement in DFD signal-to-noise over MATBG.

\subsubsection{Pentalayer Graphene on hBN}

Pentalayer graphene exhibits $T_c \approx 300$~mK with strong spin-orbit
coupling from proximity to hBN. The flat-band amplification $\mathcal{A}
\approx 80$ is lower than MATBG. The platform's advantage is cleanliness
and reproducibility, but this does not translate to DFD sensitivity.

\subsubsection{TBG/MoS$_2$ Heterostructure}

Adding a MoS$_2$ substrate to MATBG introduces spin-orbit coupling and
breaks inversion symmetry, potentially stabilising mixed-parity SC and
enhancing the P-channel interaction relevant to DFD. The $T_c$ is similar
to MATBG but the phase diagram is richer. This system is worth tracking
for future DFD parameter space exploration.

\subsubsection{La$_3$Ni$_2$O$_7$ Nickelate (Comparison Reference)}

The nickelate system (W3i3 finding, FOM$_{\rm RT} = 252$~K) has $T_c \approx
80$~K and a DoS enhancement $\mathcal{A} \approx 3$ (due to Van~Hove
singularity near the $d_{z^2}$ band). Although $\mathcal{A}$ is small, the
high $T_c$ produces $\delta T_c \approx 9$~mK --- the largest absolute DFD
$T_c$ shift of any system studied. This remains the primary DFD SC candidate.

\begin{finding}
  \label{find:system_choice}
  Among flat-band moiré systems, MATBG provides the best DFD
  signal-to-background ratio due to its combination of high $\mathcal{A}$
  and moderate $T_c$. However, no flat-band system approaches the DFD
  sensitivity of high-$T_c$ materials (nickelates, cuprates) where the
  larger gap scale amplifies the fractional DFD effect into the measurable
  regime. The optimal DFD-sensitive material remains La$_3$Ni$_2$O$_7$
  as identified in W3i3.
\end{finding}

%=============================================================================
\section{MATBG as a Psi-Sensor Element: Connection to P2 and P11}
\label{sec:sensor}
%=============================================================================

\subsection{MATBG as a Resistance-Based Psi Transducer}

Although MATBG cannot serve as structural SRF cavity material (too low
$T_c$, too small area), it can function as a \emph{psi-field sensor
element} integrated into a P2 resonator system. The operating principle:

\begin{enumerate}
  \item Embed a MATBG device near the QCP at $\nuFill \approx 2$,
    biased at $T \approx T_c/2$.
  \item The resistance $R(\psi)$ is non-linearly sensitive to $\psi_0$
    through the QCP scaling $R \propto |n - n_c(\psi_0)|^{-\nu z}$.
  \item Modulate the psi-field at frequency $f_{\rm psi}$ using a P2
    resonator (Nb$_3$Sn + YBCO SRF cavity).
  \item Measure the MATBG resistance at $f_{\rm psi}$ via lock-in
    detection.
\end{enumerate}

\subsubsection{Estimated Transducer Sensitivity}

The resistance responsivity $\partial R/\partial\psi_0$ at the QCP bias point:
\begin{align}
  \frac{\partial R}{\partial\psi_0}
    &= \frac{\partial R}{\partial n} \cdot
       \frac{\partial n}{\partial\psi_0} \nonumber\\
    &\approx \nu z \cdot \frac{R_N}{|n - n_c|} \cdot
       |\lambda-1| \cdot n_s \cdot \mathcal{A} \nonumber\\
    &\approx 1.3 \cdot \frac{26\,\text{k}\Omega}{10^9\,\text{cm}^{-2}}
       \times \lambdabound \times 6.4\ee{12} \times 143 \nonumber\\
    &\approx 3.4\ee{-7}\,\Omega/(\text{unit } \psi_0).
  \label{eq:R_responsivity}
\end{align}
With a P2 resonator driving $\psi_0 \sim 10^{-3}$ (achievable in a high-$Q$
SRF cavity, $Q \sim 10^9$, driven at 1 nW), the expected resistance signal:
$\delta R \approx 3.4\ee{-7} \times 10^{-3} \approx 3.4\ee{-10}$~$\Omega$.
This is far below the Johnson noise floor ($\delta R_J \sim
\sqrt{4k_BT R \Delta f} / I_{\rm bias} \sim 10^{-3}$~$\Omega$ at 20~mK).

\begin{finding}
  \label{find:transducer}
  MATBG as a psi-transducer element integrated with a P2 resonator
  has estimated resistance sensitivity $\partial R/\partial\psi_0 \approx
  3.4\ee{-7}$~$\Omega/\psi_0$, producing a resistance signal
  $\delta R \approx 3.4\ee{-10}$~$\Omega$ at $\psi_0 \sim 10^{-3}$.
  The Johnson noise floor exceeds this by $\sim 10^9$. The concept is
  viable in principle but requires a $10^5\times$ improvement in $|\lambda-1|$
  over current bounds, or a $10^{10}\times$ improvement in noise floor,
  to reach practical utility.
\end{finding}

\subsection{Connection to P3 (Quantum Computing)}

MATBG is directly relevant to Patent~3 (quantum control and coherence):
\begin{itemize}
  \item \textbf{Topological qubits from MATBG}: At certain $(n, D, B)$
    parameter points, MATBG may host topological SC (TSC) phases with
    Majorana zero modes (MZMs) at vortices. W3i3 showed MZMs are
    NOT lifted by DFD (particle-hole symmetry preserved), meaning
    MATBG-based topological qubits are DFD-robust --- a positive result
    for P3.
  \item \textbf{Correlated qubit gate fidelity}: DFD modifies the
    effective Hubbard-$U$ by $\delta U_{\rm DFD} \approx 2.2\ee{-6}$~meV
    (eq.~\ref{eq:deltaU_value}), producing a fractional qubit frequency
    shift $\delta\omega/\omega \approx \delta U/U \approx 2.2\ee{-7}$.
    This is $\sim 10^7$ times below current qubit frequency instability
    ($\sim 0.1\%$), confirming DFD causes negligible decoherence in
    MATBG-based qubits.
  \item \textbf{psi-field quantum error correction}: A DFD-induced
    systematic phase error of $\delta\phi \approx 2\pi \delta\omega T_{\rm gate}
    \approx 2\pi \times 2.2\ee{-7} \times 50\,\text{ns} \approx 7\ee{-14}$~rad
    per gate. This is negligible even for $10^6$ gate operations (total
    accumulated error $< 10^{-7}$~rad).
\end{itemize}

%=============================================================================
\section{Cross-Patent Integration}
\label{sec:crosspatent}
%=============================================================================

\subsection{P15 (Materials and Fabrication Methods)}

W3i4 provides the following inputs to P15:
\begin{itemize}
  \item MATBG is confirmed as a \emph{sensor-element} material class,
    not a structural DFD generator material.
  \item Required fabrication specifications:
    twist angle $\theta = 1.10^\circ \pm 0.01^\circ$;
    hBN encapsulation;
    dual graphite gate geometry;
    one-dimensional edge contacts.
  \item Operating conditions: $T < 50$~mK, $B < 0.1$~T,
    $n$ near $\nuFill = 2$ ($n_s/2 \approx 3.2\ee{12}$~cm$^{-2}$).
  \item Material ranking for DFD sensitivity (structural SRF):
    La$_3$Ni$_2$O$_7 >$ Nb$_3$Sn + YBCO laminate $>$ MATBG (sensor only).
\end{itemize}

\subsection{P2 (Resonators and Metamaterials)}

\begin{itemize}
  \item MATBG integration as a resistance-based psi-sensor in the Nb$_3$Sn
    + YBCO laminate resonator cavity (W3i3 optimal geometry).
  \item The MATBG element should be placed at the \emph{electric field
    antinode} of the SRF cavity mode, where Process~A sourcing is maximal.
  \item Required isolation: the MATBG device must be thermally isolated
    from the SRF cavity walls to prevent RF heating (thermal link
    $\kappa < 10^{-12}$~W/K required at 20~mK).
\end{itemize}

\subsection{P3 (Quantum Control and Computing)}

\begin{itemize}
  \item DFD phase shift per gate $\delta\phi < 10^{-13}$~rad: negligible.
  \item MATBG TSC phase (topological qubit platform) is DFD-robust
    (MZMs unaffected).
  \item Recommendation: MATBG is a viable topological qubit material
    for P3 applications; DFD does not compromise its quantum error
    correction threshold.
\end{itemize}

\subsection{P11 (EM Coupling and SRF Cavities)}

\begin{itemize}
  \item SRF cavity coupled to MATBG sensor element: the cavity provides
    amplified psi-field that is transduced by MATBG resistance.
  \item The Nb$_3$Sn + YBCO laminate SRF cavity (P11 optimal geometry)
    enhances the local psi-field by $\mathcal{A}_{\rm SRF} \approx 5\ee{-18}$
    (W3i2/W3i3) relative to the background --- but this is at the
    resonator's internal field, which does not exit the cavity. The
    MATBG sensor must be placed \emph{inside} the SRF cavity volume.
\end{itemize}

%=============================================================================
\section{Summary and Forward Outlook}
\label{sec:summary}
%=============================================================================

\subsection{Key Findings of W3i4}

\begin{enumerate}
  \item \textbf{MATBG flat-band enhancement}: The DoS amplification
    $\mathcal{A} \approx 143$ relative to normal graphene provides the
    largest DFD-to-phonon coupling ratio in any 2D material, but is
    still insufficient to bring the DFD signal above experimental resolution
    at the current $|\lambda-1|$ bound.

  \item \textbf{T$_c$ shift}: $\delta T_c \approx 54\,\mu$K at $|\lambda-1|
    = \lambdabound$, vs.\ experimental resolution $\sim 50$~mK.
    Gap: $\times 10^3$.

  \item \textbf{Mott gap modification}: $\delta\Vmott \approx 0.1$~neV,
    vs.\ tunnelling spectroscopy resolution $\sim 10\,\mu$eV.
    Gap: $\times 10^8$.

  \item \textbf{Phase boundary shift}: $\delta n \approx 7\ee{5}$~cm$^{-2}$,
    vs.\ gate control resolution $\sim 10^9$~cm$^{-2}$.
    Gap: $\times 10^3$.

  \item \textbf{Smoking-gun protocol}: Differential phase diagram measurement
    (shielded vs.\ unshielded EM environment) is the most discriminating
    test but currently limited by intrinsic disorder
    ($\delta n_{\rm disorder} \gg \delta n_{\rm DFD}$).

  \item \textbf{System ranking}: La$_3$Ni$_2$O$_7$ remains the optimal
    DFD-sensitive superconductor; MATBG is the best sensor element among
    flat-band systems.

  \item \textbf{Qubit safety}: DFD perturbation to MATBG-based topological
    qubits is $\sim 10^{-13}$~rad per gate, entirely negligible.

  \item \textbf{Required improvement}: Detection of DFD in MATBG would
    require either $|\lambda-1|$ to be $10^3$--$10^4\times$ larger than
    the current bound, or disorder reduction by the same factor.
\end{enumerate}

\subsection{Proposed W3i5 Directions}

\begin{itemize}
  \item \textbf{Rhombohedral graphene multilayers}: $N$-layer rhombohedral
    graphene (N = 4, 5) has flat bands without twist-angle tuning and
    potentially lower disorder; worth computing DFD FOM.
  \item \textbf{Moiré TMD at higher $T_c$}: If twisted MoS$_2$/WSe$_2$
    systems achieve $T_c > 1$~K (recent claims), recalculate DFD sensitivity.
  \item \textbf{Enhanced resonator coupling}: Model the MATBG + SRF
    cavity coupled system using full quantum electrodynamics of the
    flat-band + cavity photon system.
  \item \textbf{Spin-polarised flat bands}: At $\nuFill = 3$ (3/4 filling),
    MATBG enters a spin-polarised state; DFD coupling to magnetic order
    parameter has not yet been computed.
\end{itemize}

%=============================================================================
\section*{Equations Summary (W3i4 New Equations)}
\addcontentsline{toc}{section}{New Equations}
%=============================================================================

\noindent This iteration adds the following new equations:
\begin{itemize}
  \item Eq.~(\ref{eq:psi_eps_mu}): psi-field permittivity/permeability coupling.
  \item Eq.~(\ref{eq:VDFD}): DFD correction to electron-electron interaction.
  \item Eq.~(\ref{eq:DOS_enhancement}): flat-band DoS enhancement factor.
  \item Eq.~(\ref{eq:VDFD_flat}): flat-band amplified DFD pairing correction.
  \item Eq.~(\ref{eq:dTc}): fractional $T_c$ enhancement formula.
  \item Eq.~(\ref{eq:dTc_value}): quantitative $\delta T_c$ in MATBG.
  \item Eq.~(\ref{eq:deltaU_value}): DFD Hubbard-$U$ modification.
  \item Eq.~(\ref{eq:dMott}): fractional Mott gap change.
  \item Eq.~(\ref{eq:dn_boundary}): SC dome boundary shift formula.
  \item Eq.~(\ref{eq:dn_abs}): absolute carrier density boundary shift.
  \item Eq.~(\ref{eq:dtheta}): magic-angle shift.
  \item Eq.~(\ref{eq:dD}): displacement field shift.
  \item Eq.~(\ref{eq:R_scaling}): QCP resistance scaling.
  \item Eq.~(\ref{eq:dR_lockin}): lock-in resistance response amplitude.
  \item Eq.~(\ref{eq:R_responsivity}): MATBG psi-transducer responsivity.
\end{itemize}

Total new equations this iteration: \textbf{15}.
Total cross-patent connections: \textbf{P2, P3, P11, P15 (4 patents)}.

\end{document}
